\documentclass[11pt]{article}
\usepackage[margin=1in]{geometry}
\usepackage{amsmath,amssymb}
\usepackage[hidelinks]{hyperref}

\newcommand{\C}{\mathbb C}
\newcommand{\Max}{\operatorname{Max}}

\begin{document}

\section*{Human Review: Square-zero unitization and non-open multiplication}

\noindent Original automatic proof: \texttt{solution\_packet.pdf}

\noindent Checked date: 18 June 2026

\noindent Status: Manually checked.

\noindent Verdict: The proof appears correct.

\section*{Human Comments}

The proposed example appears to correctly answer the first clause of Open
Problem 3 in Draga--Kania as printed.  Namely, it gives a commutative unital
Banach algebra whose maximal ideal space is zero-dimensional, but whose
multiplication map is not open.

The construction is the unitization of a square-zero ideal.  Let
\[
  A=\C\oplus E,
  \qquad E=\C^2,
\]
with multiplication
\[
  (\lambda,x)(\mu,y)=(\lambda\mu,\lambda y+\mu x)
\]
and norm \(\|(\lambda,x)\|=|\lambda|+\|x\|_\infty\).  The algebraic and norm
checks are straightforward: \(A\) is a commutative unital Banach algebra, and
the norm is submultiplicative.

\section*{Maximal Ideal Space}

The computation of the maximal ideal space is correct.  If
\(\varphi:A\to\C\) is a character and \(x\in E\), then
\[
  (0,x)^2=0.
\]
Therefore
\[
  \varphi(0,x)^2=\varphi((0,x)^2)=0,
\]
so \(\varphi(0,x)=0\).  Since any character on a unital complex Banach
algebra satisfies \(\varphi(1)=1\), it follows that
\[
  \varphi(\lambda,x)=\lambda.
\]
Thus \(A\) has exactly one character, so \(\Max(A)\) is a singleton and hence
zero-dimensional.

\section*{Failure of Openness}

The failure of openness is also correct.  Let \(e_1=(1,0)\), \(e_2=(0,1)\) in
\(E\), and put
\[
  a=(0,e_1).
\]
Then \(a^2=0\).  To show that multiplication is not open at \((a,a)\), it is
enough to find \(0<\varepsilon<1\) such that
\[
  B_A(a,\varepsilon)\cdot B_A(a,\varepsilon)
\]
is not a neighbourhood of \(0\).

Given any \(\delta>0\), choose \(0<t<\delta\) and consider
\[
  z=(0,t e_2).
\]
Then \(\|z\|<\delta\).  Suppose, for contradiction, that
\[
  z=(\alpha,x)(\beta,y)
\]
with both factors in \(B_A(a,\varepsilon)\).  Comparing scalar coordinates
gives
\[
  \alpha\beta=0.
\]
Hence \(\alpha=0\) or \(\beta=0\).  If \(\alpha=0\), then the ideal
coordinate equation becomes
\[
  \beta x=t e_2.
\]
Since \(t\neq0\), one has \(\beta\neq0\), and so \(x\in \C e_2\).  But every
vector in \(\C e_2\) is at distance at least \(1\) from \(e_1\) in the
supremum norm:
\[
  \|x-e_1\|_\infty\geq 1.
\]
This contradicts the assumption that
\[
  (\alpha,x)\in B_A(a,\varepsilon),
  \qquad \varepsilon<1.
\]
The case \(\beta=0\) is identical, using \(y\) instead of \(x\).

Thus every ball around \(0=a^2\) contains points \((0,t e_2)\) which are not
in the product of the two \(\varepsilon\)-balls around \(a\).  Therefore
multiplication is not open.

\section*{Review Outcome}

The proof appears correct.  It is a simple and effective counterexample to
the first clause of the open problem as literally stated.

The main caveat is semantic rather than mathematical: the example is highly
non-semisimple, and it has only the trivial idempotents.  Thus it does not
answer the stronger dense-idempotent follow-up, and it may be worth asking the
authors whether they intended an additional assumption such as semisimplicity.
As written, however, the counterexample is valid.

\end{document}
